导数观察局部变化，会把高频噪声和舍入误差放大；积分汇总区间贡献，能平均局部扰动。两类方法都可从"用低次多项式替代函数"推导，但误差传播完全不同。

以中心差商近似 \(f'(x)\) 为例。取 \(f(x)=\sin x\)，\(x=\pi/4\)，真值为 \(\cos(\pi/4)\approx 0.70710678\)。用步长 \(h=10^{-3}\)，差分得 \(0.70710666\)，约 8 位正确。把 \(h\) 缩小到 \(10^{-8}\)，截断误差可降至 \(10^{-16}\) 量级——但浮点舍入把 \((f(x+h)-f(x-h))/(2h)\) 的分子推向了消去区，有效数字反而只剩约 8 位。步长不是越细越好。

本单元从步长的 U 形误差开始，再把梯形、Simpson、Gaussian 与自适应求积连成一条选择路径。

\begin{enumerate}
\tightlist
\item
  数值微分为何不能无限缩小步长
\item
  从插值多项式到梯形与 Simpson 求积
\item
  Gaussian 与自适应求积
\end{enumerate}

同一个 \(h\) 扮演两种角色：变小降低截断误差，却也可能放大舍入、噪声和采样盲区。公式的阶数只描述进入渐近区后的理想行为；实际选择还要看函数是否光滑、能否自由取样，以及粗细两套节点是否可能同时漏掉结构。

前一单元：\hyperref[ch:063]{``插值与逼近：通过数据点，不等于理解点之间''}。

\subsection{本章篇目}

以下篇目按概念依赖排列；若从具体问题进入，也可以先打开最接近当前困惑的一篇，再沿篇内链接回补。

\begin{itemize}
\tightlist
\item \hyperref[sec:08-Numerical-Analysis/04-Numerical-Differentiation-and-Integration/01]{``数值微分为何不能无限缩小步长''}；
\item \hyperref[sec:08-Numerical-Analysis/04-Numerical-Differentiation-and-Integration/02]{``从插值多项式到梯形与 Simpson 求积''}；
\item \hyperref[sec:08-Numerical-Analysis/04-Numerical-Differentiation-and-Integration/03]{``Gaussian 与自适应求积''}。
\end{itemize}
